\chapter{総合演習}\label{章:総合演習}

ここでは,問題のすぐ前に使う公式を置かず,式を見て方針を選ぶ練習をする.計算を始める前に,求めるものが原始関数,定積分の値,関数の性質,極限,あるいは誤差付きの近似のどれであるかを確かめたい.同じ積分記号でも,それによって使いやすい手法は変わる.

基礎問題は1つか2つの基本操作を確実につなぐ問題,標準問題は積分区間の分け方や式の変形を選ぶ問題である.発展問題では複数の章の内容を使うが,新しい一般定理を調べることは求めない.微分した式を元へ戻す,平方の積分を考える,端点の極限を独立に調べるなど,本書で扱った操作から進められる.ヒントには各問題を何に帰着させるかと最初の変形を書き,問題の後にまとめた.解答は巻末に置いた.最初の一手を自分で考えてから,必要な番号だけを参照してほしい.

\section{基礎問題}\label{節:総合基礎}

\begin{enumerate}[label=\textbf{\thechapter.\arabic*},leftmargin=*]
 \item\label{問:総合一割り算}
 次の定積分を求めよ.
 \[
  \int_0^1\frac{x^3+x^2+2x+1}{1+x^2}\dd x.
 \]
 \item\label{問:総合二置換}
 次の定積分を求め,置換後の区間も示せ.
 \[
  \int_0^1x\sqrt{1+x^2}\dd x.
 \]
 \item\label{問:総合三部分積分}
 次の定積分を求めよ.
 \[
  \int_1^{\e}x\log x\dd x.
 \]
 \item\label{問:総合四動く端点}
 $x>0$に対して$F(x)=\int_x^{x^2}\sqrt{1+t^2}\dd t$とする.
 $F'(x)$を求め,曲線$y=F(x)$の$x=1$における接線を求めよ.
 \item\label{問:総合五相補重み}
 $a>0$とする.次の積分を$a$で表せ.
 \[
  \int_{-a}^{a}\frac{1+x^2}{1+\e^x}\dd x.
 \]
 \item\label{問:総合六面積体積}
 曲線$y=x^2$と直線$y=x$で囲まれる領域の面積を求めよ.
 また,この領域を$x$軸のまわりに1回転してできる立体の体積を求めよ.
\end{enumerate}


%
\section{標準問題}\label{節:総合標準}

\begin{enumerate}[label=\textbf{\thechapter.\arabic*},leftmargin=*,start=7]
 \item\label{問:総合七区間の最適化}
 実数$a$に対して
 \[
  H(a)=\int_a^{a+1}(x^2-2x)^2\dd x
 \]
 とする.\,$H(a)$の最小値と,それを与える$a$をすべて求めよ.
 \item\label{問:総合八置換と部分積分}
 次の定積分を求めよ.
 \[
  \int_0^1x^3\log(1+x^2)\dd x.
 \]
 \item\label{問:総合九区分求積}
 次の極限を求めよ.
 \[
  \lim_{n\to\infty}\sum_{k=1}^{n}\frac{k^2}{n(n^2+k^2)}.
 \]
 \item\label{問:総合十軸を横切る回転}
 $0\leqq x\leqq1$,\,$x-1\leqq y\leqq x$で定まる領域を$x$軸のまわりに一回転してできる立体の体積を求めよ.
 各$x$での回転断面が円板になるか円環になるかを説明し,半径が切り替わる点を求めること.
 \item\label{問:総合十一漸化式}
 非負整数$n$に対して$K_n=\int_0^1\frac{x^n}{1+x}\dd x$とする.
 $K_0$と,\,$n\geqq1$で成り立つ漸化式を求め,\,$K_3$を計算せよ.
 さらに$K_n\to0$を示し,それを使って
 \[
  \lim_{n\to\infty}\sum_{k=1}^{n}\frac{(-1)^{k-1}}{k}=\log2
 \]
 を示せ.
 \item\label{問:総合十二両端判定}
 実数$\alpha$に対して,広義積分
 \[
  \int_0^{\infty}\frac{x^\alpha}{1+x^2}\dd x
 \]
 の収束条件を求めよ.また$\alpha=0$の場合の値を求めよ.
\end{enumerate}


%
\section{発展問題}\label{節:総合発展}

\begin{enumerate}[label=\textbf{\thechapter.\arabic*},leftmargin=*,start=13]
 \item\label{問:総合十三積分方程式}
 $[0,1]$で連続な関数$f$が,すべての$x\in[0,1]$について
 \[
  f(x)=1+\int_0^x(x-t)f(t)\dd t
 \]
 を満たすとする.\,$f$を求め,求めた関数が確かにこの等式を満たすことを確かめよ.
 微分方程式の一般的な解法は用いず,積の微分と基本定理から求めてよい.
 \item\label{問:総合十四二つの制約}
 $[0,1]$で連続な実数値関数$f$が
 \[
  \int_0^1f(x)\dd x=1,\qquad
  \int_0^1xf(x)\dd x=0
 \]
 を満たすとする.\,$\int_0^1f(x)^2\dd x$の最小値と,それを与える関数を求めよ.
 \item\label{問:総合十五媒介曲線}
 $0\le t\le\pi$で
 \[
  x=t-\frac12\sin2t,\qquad y=\sin^2t
 \]
 と表される曲線を考える.この曲線と$x$軸で囲まれる領域の面積,曲線の長さ,領域を$x$軸のまわりに一回転してできる立体の体積を求めよ.
 面積と体積の積分で,同じ部分を重複して数えていないことも確かめること.
 \item\label{問:総合十六近似と条件収束}
 次の2つの問いに,原始関数の初等的な表示を求めずに答えよ.
 \begin{enumerate}[label=(\arabic*),leftmargin=*]
  \item 有限テイラー公式を使い,\,$A=\int_0^1\sin(x^2)\dd x$を幅$\frac1{10000}$以下の有理数の区間に挟め.
  \item 広義積分$\int_1^{\infty}\sin(x^2)\dd x$が条件収束することを示せ.
 \end{enumerate}
 \item\label{問:総合十七ウォリス}
 非負整数$n$に対して
 \[
  C_n=\int_0^1(1-x^2)^n\dd x
 \]
 とする.\,$C_0$と$C_n$の漸化式を求め,\,$C_n$を有限の積で表せ.
 次にウォリスの積分との関係を使って,
 \[
  \lim_{n\to\infty}\sqrt n\,C_n,
  \qquad
  \lim_{n\to\infty}n^{\frac32}(C_n-C_{n+1})
 \]
 を求めよ.
 \item\label{問:総合十八相補と半角}
 $n$を正の整数とする.次の広義積分を,\,$n$を含む有限の積で表せ.
 \[
  J_n=\int_{-\pi}^{\pi}
  \frac{\sin^{2n}x}{(1+\e^x)\sqrt{1+\cos x}}\dd x.
 \]
 最初に両端での収束を確かめること.対称性,半角公式,置換積分,部分積分を組み合わせれば,ベータ関数やガンマ関数の公式を使わずに求められる.
\end{enumerate}

\clearpage

%
\section{融合問題}\label{節:融合問題}
以下では,数列・図形・整数・組合せ・ベクトル・確率に関係する問題を扱う.新しい大学数学の定理は使わず,必要な考え方は問題文と解答で説明する.
\begin{enumerate}[label=\textbf{\thechapter.\arabic*},leftmargin=*,itemsep=6pt,start=19]
\item\label{問:融合ニュートン}
$x_1=2$,\,$x_{n+1}=\frac12\left(x_n+\frac2{x_n}\right)$とする.
\begin{enumerate}[label=(\arabic*),leftmargin=*]
\item この漸化式が,方程式$x^2-2=0$に対するニュートン法から得られることを示せ.
\item $b_n=\frac{x_n-\sqrt2}{x_n+\sqrt2}$と置き,\,$b_{n+1}$を$b_n$で表せ.これを用いて$x_n$を$n$の式で表せ.
\item $q=\frac{2-\sqrt2}{2+\sqrt2}$として,\,$\lim_{n\to\infty}\frac{x_n-\sqrt2}{q^{2^{n-1}}}$を求めよ.
\end{enumerate}
\item\label{問:融合放物線}
放物線$y=x^2$上の2点$A(a,a^2),B(b,b^2)$を考える.ただし$a<b$とする.両点における接線の交点を$T$,放物線上で$x$座標が$\frac{a+b}{2}$である点を$C$とする.
\begin{enumerate}[label=(\arabic*),leftmargin=*]
\item $T$の座標を求めよ.また$AB$に平行な接線の接点が$C$であることを示せ.
\item 弦$AB$と放物線に囲まれた面積$Q$,三角形$ABC$の面積$P$,三角形$ABT$の面積$R$を求め,比$Q:P:R$を求めよ.
\item $a+b=2c$を固定するとき,\,$T$の軌跡を求めよ.
\end{enumerate}
\item\label{問:融合自然対数の底}
正の整数$n$について,\,$f_n(x)=\frac{x^n(1-x)^n}{n!}$,\,$I_n=\int_0^1\e^x f_n(x)\dd x$とする.
\begin{enumerate}[label=(\arabic*),leftmargin=*]
\item $0\leqq k\leqq2n$のすべての整数$k$について,\,$f_n^{(k)}(0),f_n^{(k)}(1)$が整数であることを示せ.ただし$f_n^{(0)}=f_n$とする.
\item $I_n=A_n\e+B_n$を満たす整数$A_n,B_n$が存在することを示せ.
\item $0<I_n\leqq\frac{\e}{4^n n!}$を示し,\,$\e$が無理数であることを証明せよ.
\end{enumerate}
\item\label{問:融合二項係数}
正の整数$m$と非負整数$n$について
\[
 S_{m,n}=\sum_{k=0}^n\frac{(-1)^k\binom nk}{m+k}
\]
とする.
\begin{enumerate}[label=(\arabic*),leftmargin=*]
\item $S_{m,n}$を1つの定積分で表し,階乗を用いて値を求めよ.
\item $m$を固定するとき,\,$\lim_{n\to\infty}n^m S_{m,n}$を求めよ.
\end{enumerate}
\item\label{問:融合近似とベクトル}
実数$a,b$に対し,\,$E(a,b)=\int_0^1\{x^2-(ax+b)\}^2\dd x$とする.
\begin{enumerate}[label=(\arabic*),leftmargin=*]
\item $p(x)=x^2-x+\frac16$について,\,$\int_0^1p(x)\dd x=\int_0^1xp(x)\dd x=0$を示せ.
\item $E(a,b)$の最小値と,それを与える$a,b$を求めよ.
\item 原点から点$(u_1,u_2)$へ向かう矢印をベクトル$\boldsymbol u=(u_1,u_2)$で表す.その長さは$|\boldsymbol u|=\sqrt{u_1^2+u_2^2}$である.2つのベクトル$\boldsymbol u=(u_1,u_2)$,\,$\boldsymbol v=(v_1,v_2)$の\textbf{内積}を$\boldsymbol u\cdot\boldsymbol v=u_1v_1+u_2v_2$とする.等式
\[
 |\boldsymbol u+\boldsymbol v|^2=|\boldsymbol u|^2+2\boldsymbol u\cdot\boldsymbol v+|\boldsymbol v|^2
\]
と,(2)の計算を比較せよ.積分のどの式が内積の役割をしているか説明せよ.
\end{enumerate}
\Needspace{17\baselineskip}
\item\label{問:融合確率と距離}
一辺$1$の正方形$0\leqq x\leqq1,0\leqq y\leqq1$から点$P(x,y)$を無作為に選ぶ.ある部分に入る確率は,その部分の面積に等しいものとする.原点から$P$までの距離を$R$とする.
\begin{enumerate}[label=(\arabic*),leftmargin=*]
\item $0\leqq r\leqq\sqrt2$について,\,$R\leqq r$となる確率を求めよ.逆三角関数を使う場合は,角の範囲も明記すること.
\item 距離の平均値は,次のように表される.
\[
 M=\int_0^1\left\{\int_0^1\sqrt{x^2+y^2}\dd y\right\}\dd x.
\]
内側を先に積分し,その結果を$x$について積分する記号である.正方形を対角線で分けることにより,
\[
 M=2\int_0^1\left\{\int_0^x\sqrt{x^2+y^2}\dd y\right\}\dd x
\]
と書けることを説明し,\,$M$を求めよ.
\end{enumerate}
\end{enumerate}
%
\clearpage
\section{ヒント}\label{節:総合ヒント}

以下では,計算の結論ではなく,最初の変形や見るべき条件を示す.解き直すときには,その変形を選んだ理由も言葉にしてみよう.

\begin{enumerate}[label=\textbf{\thechapter.\arabic*},leftmargin=*]
 \item 分子を$1+x^2$で割ると,余りの項が分母の微分と結び付く.
 \item 根号の中を$u=1+x^2$と置く.\,$x\dd x$をそのまま$\dd u$で表す.
 \item $\log x$を微分する側に選ぶと,残る積分は整式になる.
 \item $G(u)=\int_0^u\sqrt{1+t^2}\dd t$と置いて,\,$F(x)=G(x^2)-G(x)$と書く.
 \item $x=-t$で書き直した積分と元の積分を足す.\,$1+x^2$は偶関数である.
 \item 交点を求めた後で,上下関係と,回転軸からの二つの距離を調べる.
 \item 端点を微分すると$H'(a)$が得られる.値の計算には$y=a-\frac12$と置き,積分区間の中心から測ると整理しやすい.
 \item $u=1+x^2$と置けば,\,$x^3\dd x=\frac12(u-1)\dd u$になる.その後で部分積分する.
 \item 一項を$\frac1n$と$\frac{k}{n}$だけを使って表す.積分に直したら,分子を分母と定数の差にする.
 \item 区間$[x-1,x]$は常に$0$を含む.回転後の半径は,両端の軸からの距離の大きい方である.
 \item $K_n+K_{n-1}$では分母が消える.極限には$0\leqq K_n\leqq\int_0^1x^n\dd x$を使う.
 \item $0$付近では$x^\alpha$,無限大では$x^{\alpha-2}$と比較し,二つの条件を同時に満たす範囲を求める.
 \item 積分を$x\int_0^xf(t)\dd t-\int_0^xtf(t)\dd t$に分けて微分する.\,$f''=f$が得られたら,\,$\e^{-x}(f+f')$の微分を考える.
 \item まず二つの制約を満たす一次関数$\ell(x)$を求める.\,$g=f-\ell$と置いたとき,\,$\int_0^1\ell g\dd x$がどうなるかを調べる.
 \item $x'(t)=2\sin^2t$を使う.弧長では$\sqrt{x'(t)^2+y'(t)^2}$を整理してから積分する.
 \item (1) $\sin u$の六次テイラー多項式では,偶数次の係数が$0$になる.七階微分の符号と大きさを調べる.(2) $u=x^2$と置く.収束には部分積分,絶対値の発散には長さ$\pi$の区間ごとの下からの評価を使う.
 \item $x=\sin t$によって$C_n$を$\cos$の奇数べきの積分にする.隣り合うウォリス積分の比と積を組み合わせる.
 \item 指数関数を含む重みは,\,$x$を$-x$にすると相補的になる.その後,\,$x=2t$,さらに$u=\sin t$と置く.残る積分の$(1-u^2)$の指数は整数なので,部分積分で一つずつ下げられる.

 \item 分子と分母をそれぞれ平方に直し,\,$b_{n+1}=b_n^2$を確かめる.
 \item 弦と放物線の差は$(x-a)(b-x)$となる.三角形は,弦の中点を通る縦の線分で二つに分ける.
 \item 二項定理で展開し,端点で残る導関数を調べる.部分積分では$\e^x$を積分し続ける.
 \item $\frac1{m+k}$を$x^{m+k-1}$の積分で表し,二項定理を使う.
 \item $x^2-(ax+b)=p(x)+(1-a)x-\frac16-b$と書き,平方を展開する.
 \item (1) $r=1$を境に,四分円が正方形からはみ出す.(2) 下側の三角形で$y=xt$と置く.内側の積分では$x$を固定する.
\end{enumerate}

\smallskip\noindent 解答は\pageref{解答:第13章}ページ.

%
